% !PRINTED 623
\weberpage{623}
If we form from (2) the $p$th power of $\omega$, and observe
\gesperrt{Fermat's} theorem for rational numbers $[a_0^p \equiv a_0 \pmod{p}]$,
we obtain
\[
\omega^p \equiv a_0 + a_1 r^p + a_2 r^{2p} + \cdots + a_{\mu-1} r^{(\mu-1)p} \pmod{p},
\]
or by (2)
\weq{3}{\omega^p \equiv \omega_p \pmod{p}.}

This determines the group $\Psi$ to which the ideal $\mathfrak{p}$
belongs. For by (3) we have also
\[
\omega^p \equiv \omega_p \pmod{\mathfrak{p}},
\]
hence $\psi_0 = (r, r^p)$, and by \S.\,162, (13) the group $\Psi$ consists
of the powers of this substitution, so far as they are different
from one another.

If therefore $f$ is the least positive exponent for which
\weq{4}{p^f \equiv 1 \pmod{m}}
holds, that is, if $p$ belongs to the exponent $f$ for the modulus $m$, then
\weq{5}{\Psi = (r,r), (r,r^p), (r,r^{p^2}), \ldots (r,r^{p^{f-1}})}
and it is of degree $f$. Accordingly $\mathfrak{p}$ is also a prime ideal
of degree $f$, and
\weq{6}{N(\mathfrak{p}) = p^f.}

The number $e$ of the conjugate prime factors of $f$ that are different
from one another is $\mu : f$, and we therefore have the theorem: % [sic] source prints $f$ after "of"; mathematically $p$ is required

III. \gesperrt{If $p$ is a prime different from $q$ which belongs to the exponent $f$ for the modulus $m$, and if further $\mu = \varphi(m) = ef$, then $p$ splits in $\Omega_m$ into $e$ conjugate prime factors of degree $f$, different from one another.}

\begin{center}
\S.\,172.

\gesperrt{The conjugate prime ideals.}
\end{center}

We must enter somewhat more closely into the formation of the conjugate prime factors $\mathfrak{p}_n$ of a prime $p$ different from $q$,
which by the theorem last proved splits in $\Omega_m$ into $e$
prime factors different from one another, when
\weq{1}{\varphi(m) = ef}
holds, and $f$ is the least positive exponent for which
\weq{2}{p^f \equiv 1 \pmod{m}.}
